\begin{proof}
At the trial point $t=q$, the denominator in
\eqref{I:F:eq:f2-scalar-reduction} equals $q^2$ exactly.  Lemma~
\ref{I:F:lem:f2-h-expansion} therefore gives the lower bound
\begin{equation}
 \mathcal P(q)\ge\frac{h(q)}q
              =\frac14+\frac14q+O(q^2).
 \label{I:F:eq:f2-P-trial}
\end{equation}
We next localize every global scalar maximizer.  For $0\le q,t\le1$,
\[
 1-\sqrt{(1-q)(1-t)}
 \ge\frac{q+t-qt}{2}\ge\frac{q+t}{4}.
\]
Together with $h(t)\le t$, this gives, for $x=t/q$,
\begin{equation}
 \frac{q h(t)}{(1-\sqrt{(1-q)(1-t)})^2}
 \le\frac{16qt}{(q+t)^2}
 =\frac{16x}{(1+x)^2}.
 \label{I:F:eq:f2-scalar-coarse-localization}
\end{equation}
For $x\notin I=[1/100,100]$ the last expression is less than $1/5$.
The trial value \eqref{I:F:eq:f2-P-trial} exceeds $1/5$ for all sufficiently
small $q$.  Thus every global maximizing point has $x\in I$.
In particular $t=qx\to0$ uniformly on all possible maximizers, which
permits the global envelope expansion to be used uniformly.

On the fixed compact interval $I$, Taylor expansion gives
\begin{align}
 1-\sqrt{(1-q)(1-qx)}
 &=\frac q2(1+x)+\frac{q^2}{8}(x-1)^2+O_I(q^3),\notag\\
 \frac{q h(qx)}{(1-\sqrt{(1-q)(1-qx)})^2}
 &=f_0(x)+qf_1(x)+O_I(q^2),
 \label{I:F:eq:f2-uniform-scalar-expansion}
\end{align}
where
\begin{equation}
 f_0(x)=\frac{x}{(1+x)^2},\qquad
 f_1(x)=\frac{x}{(1+x)^2}
       \left(x-\frac{(x-1)^2}{2(1+x)}\right).
 \label{I:F:eq:f2-f0-f1}
\end{equation}
The remainder is uniform in $x\in I$: the numerator expansion comes
from \eqref{I:F:eq:f2-h-expansion}, and the leading denominator is bounded
away from zero after division by $q^2$.

The exact identity
\[
 \frac14-f_0(x)=\frac{(x-1)^2}{4(1+x)^2}
\]
gives a quadratic maximum at $x=1$.  On $I$ there is $c_I>0$ such that
$f_0(x)\le1/4-c_I(x-1)^2$.  Also $f_1(1)=1/4$ and
$|f_1(x)-f_1(1)|\le L_I|x-1|$ for a finite Lipschitz constant $L_I$.
Completing the square yields uniformly on $I$
\begin{align*}
 f_0(x)+qf_1(x)
 &\le\frac14+\frac q4-c_I(x-1)^2+qL_I|x-1|\\
 &\le\frac14+\frac q4+\frac{L_I^2}{4c_I}q^2.
\end{align*}
Combining this with \eqref{I:F:eq:f2-uniform-scalar-expansion} gives the
matching upper bound.  Equation \eqref{I:F:eq:f2-P-trial} then proves
\eqref{I:F:eq:f2-P-expansion}.
\end{proof}
